\documentclass[11pt]{article}
\usepackage{amsmath,amsthm,amssymb}
\usepackage[margin=1in]{geometry}
\usepackage{hyperref}
\usepackage{enumitem}

\newtheorem{theorem}{Theorem}
\newtheorem{lemma}[theorem]{Lemma}
\newtheorem{corollary}[theorem]{Corollary}
\newtheorem{proposition}[theorem]{Proposition}
\newtheorem{conjecture}[theorem]{Conjecture}
\theoremstyle{definition}
\newtheorem{remark}[theorem]{Remark}
\newtheorem{definition}[theorem]{Definition}
\newtheorem{example}[theorem]{Example}

\newcommand{\Hq}{H_q}
\newcommand{\Sn}{S_n}
\newcommand{\Vlam}{V^{\lambda}}
\newcommand{\Om}{\Omega}
\newcommand{\hatl}{\widehat\lambda}
\newcommand{\elh}{\ell(\hatl)}
\newcommand{\Rp}{{R'}}
\newcommand{\SYT}{\mathrm{SYT}}
\newcommand{\tr}{\mathrm{tr}}
\newcommand{\qint}[1]{[#1]_q}
\newcommand{\Trs}{T_{\mathrm{rs}}}
\newcommand{\Tstar}{T^*}

\title{Per-SYT positivity of $\Om^{(\lambda)}$ in the asymmetric Hoefsmit
basis, and the converse trace-vanishing direction}
\author{Clio}
\date{17 May 2026}

\begin{document}
\maketitle

\begin{abstract}
Let $\hatl$ be a partition with $\elh \ge 2$ and every row of length
$\ge 2$, and let $\lambda = (\hatl, 1^q) \vdash n$.  Working in the
asymmetric Hoefsmit seminormal basis with off-diagonal normalization
$b' = 1$, we prove that every matrix entry of $\Om^{(\lambda)}$ on
$\Vlam$ is a non-negative rational function of~$q$, i.e.\ belongs to
$\mathbb{Q}_{\ge 0}(q) := \{N(q)/D(q) : N, D \in \mathbb{Z}_{\ge 0}[q]\}$
in lowest terms.  This is the structural input for the converse
trace-vanishing direction: if $\tau(\hatl) = 0$ then
$\tr_{\Vlam}(\Om^{(\lambda)}) \ne 0$, modulo the existence of a
descent-$0$-at-$1$ SYT of $\lambda$ with no column-descents
(Conjecture~\ref{conj:T-exists}).  We give an explicit construction
covering all cases tested for $|\lambda| \le 14$, in particular $q = 1$
and a wide regime of $q \ge 2$ shapes.
\end{abstract}

\section{Setup}\label{sec:setup}

We follow the conventions of \cite{multirowMay15} and
\cite{traceVanishingMay16}.  $\Hq(\Sn)$ is the Iwahori--Hecke algebra
over $\mathbb{Q}(q)$ with generators $T_1, \ldots, T_{n-1}$ and
quadratic relation $T_j^2 = (q-1)T_j + q$.  Set $S_j := T_j + 1$ and
\[
   \Rp_k \;:=\; S_{k-1} S_{k-2} \cdots S_1, \qquad
   \Om^{(\nu)} \;:=\; \Rp_{\ell(\nu)+1}^{(\nu)} \cdots \Rp_{|\nu|}^{(\nu)}.
\]

Throughout, $\hatl$ is a partition with $\elh \ge 2$ and every row
$\ge 2$, and $\lambda = (\hatl_1, \ldots, \hatl_{\elh}, 1^q) \vdash n$,
$n = |\hatl| + q$, $\ell_\lambda = \elh + q$, with
\[
   \tau(\hatl) := \max\bigl(0,\ q + 2\elh - 1 - |\hatl|\bigr).
\]
$\Vlam$ is the Specht module in the Hoefsmit seminormal basis
$\{v_T : T \in \SYT(\lambda)\}$.

\paragraph{The asymmetric Hoefsmit ($b' = 1$) basis.}
For each $T \in \SYT(\lambda)$ and each $i \in \{1, \ldots, n-1\}$, let
$d_i(T) := c_{i+1}(T) - c_i(T)$ be the axial distance from $i$ to $i+1$,
where $c(r, k) = k - r$ is the content of cell $(r,k)$.  The action of
$T_i$ on $v_T$ is:
\begin{itemize}[topsep=2pt,itemsep=2pt]
\item $d = 1$ (same row): $T_i v_T = q\, v_T$.
\item $d = -1$ (same column): $T_i v_T = -v_T$.
\item $|d| \ge 2$ (block case, with $s_i T \in \SYT(\lambda)$): in the
asymmetric ($b' = 1$) convention,
\begin{align}
   T_i v_T &= a_d\, v_T + 1 \cdot v_{s_i T},
\label{eq:Ti-T}\\
   T_i v_{s_i T} &= (a_d a_{-d} + q)\, v_T + a_{-d}\, v_{s_i T},
\label{eq:Ti-sT}
\end{align}
where
\[
   a_d := \frac{(q-1)\,q^d}{q^d - 1}.
\]
\end{itemize}

\section{The positivity lemma}\label{sec:lemma1}

\begin{lemma}[Hoefsmit positivity]\label{lem:positivity}
Every entry of the matrix of $S_i = T_i + 1$ on $\Vlam$ in the
asymmetric Hoefsmit basis lies in $\mathbb{Q}_{\ge 0}(q)$.  Explicitly:
\begin{enumerate}[label=(\roman*),topsep=2pt,itemsep=2pt]
\item Diagonal entries:
\[
   [S_i]_{T,T} \;=\;
   \begin{cases}
      q+1, & d_i(T) = 1,\\
      0,   & d_i(T) = -1,\\
      \dfrac{\qint{d+1}}{\qint{d}}, & |d_i(T)| \ge 2,
   \end{cases}
\]
where $\qint{k} = (q^k - 1)/(q-1)$.  Each non-zero diagonal value
belongs to $\mathbb{Z}_{\ge 0}[q] / \mathbb{Z}_{\ge 0}[q]$ in lowest
terms.
\item Off-diagonal entries (for $|d| \ge 2$):
\[
   [S_i]_{s_i T, T} = 1, \qquad
   [S_i]_{T, s_i T} \;=\;
   q\,\frac{\qint{d-1}\,\qint{d+1}}{\qint{d}^{\,2}}.
\]
Both belong to $\mathbb{Q}_{\ge 0}(q)$.  All other off-diagonal
entries of $S_i$ are zero.
\end{enumerate}
\end{lemma}

\begin{proof}
The action $T_i v_T = q\,v_T$ for $d = 1$ gives $S_i v_T = (q+1) v_T$;
the action $T_i v_T = -v_T$ for $d = -1$ gives $S_i v_T = 0$.

For $|d| \ge 2$, from \eqref{eq:Ti-T}--\eqref{eq:Ti-sT}:
\[
   [S_i]_{T,T} = a_d + 1 = \frac{(q-1)q^d + (q^d - 1)}{q^d - 1}
              = \frac{q^{d+1} - 1}{q^d - 1} = \frac{\qint{d+1}}{\qint{d}}.
\]
For $d \ge 2$, both $\qint{d+1}$ and $\qint{d}$ are polynomials in $q$
with positive integer coefficients, so the ratio is in
$\mathbb{Z}_{\ge 0}[q]/\mathbb{Z}_{\ge 0}[q]$.  For $d \le -2$, set
$d = -k$ with $k \ge 2$; using $\qint{-m} = -q^{-m}\,\qint{m}$,
\[
   \frac{\qint{-k+1}}{\qint{-k}}
   = \frac{-q^{-(k-1)}\,\qint{k-1}}{-q^{-k}\,\qint{k}}
   = q\,\frac{\qint{k-1}}{\qint{k}}
   \;\in\; \mathbb{Q}_{\ge 0}(q).
\]

For the off-diagonal $[S_i]_{s_i T, T}$, equation \eqref{eq:Ti-T} gives
$[T_i]_{s_i T, T} = 1$, hence $[S_i]_{s_i T, T} = 1 \in \mathbb{Z}_{\ge 0}[q]$.

For $[S_i]_{T, s_i T}$, equation \eqref{eq:Ti-sT} gives
$[T_i]_{T, s_i T} = a_d a_{-d} + q$.  Computing,
\[
   a_d a_{-d} = \frac{(q-1)^2 q^d}{(q^d - 1)(1 - q^d)}
              = -\frac{(q-1)^2 q^d}{(q^d-1)^2},
\]
so
\[
   a_d a_{-d} + q
   = \frac{q (q^d - 1)^2 - (q-1)^2 q^d}{(q^d-1)^2}
   = \frac{(q^{d+1} - q^2)(q^d - q^{-1})}{(q^d - 1)^2}
   = \frac{q(q^{d-1} - 1)(q^{d+1} - 1)}{(q^d-1)^2}.
\]
Rewriting numerator and denominator in $\qint{\cdot}$:
\[
   [S_i]_{T, s_i T} \;=\; q\,\frac{\qint{d-1}\,\qint{d+1}}{\qint{d}^{\,2}}.
\]
For $d \ge 2$ this is manifestly a quotient of non-negative-coefficient
polynomials.  For $d = -k \le -2$, using
$\qint{-k\pm 1} = -q^{-(k\mp 1)}\qint{k\mp 1}$,
\[
   q\,\frac{\qint{-k-1}\,\qint{-k+1}}{\qint{-k}^{\,2}}
   = q\,\frac{\qint{k+1}\,\qint{k-1}}{\qint{k}^{\,2}}
   \;\in\; \mathbb{Q}_{\ge 0}(q).
\]
All other entries of $S_i$ are zero by definition of the basis.
\end{proof}

\begin{remark}
Lemma~\ref{lem:positivity} can be verified symbolically for $|d|$ up to
any chosen bound; it has been checked computationally for
$|d| \le 6$ in
\texttt{\textasciitilde/projects/scratch/2026-05-17-converse-prove/verify\_lemma1.py}.
\end{remark}

\begin{corollary}[Product positivity]\label{cor:product-positivity}
For any product of $S_i$'s, the matrix on $\Vlam$ in the asymmetric
Hoefsmit basis has every entry in $\mathbb{Q}_{\ge 0}(q)$.  In
particular, every entry of $\Om^{(\lambda)}$ on $\Vlam$ is in
$\mathbb{Q}_{\ge 0}(q)$.
\end{corollary}

\begin{proof}
The set $\mathbb{Q}_{\ge 0}(q)$ is closed under sums and products: if
$f = N_1/D_1$ and $g = N_2/D_2$ with $N_i, D_i \in \mathbb{Z}_{\ge 0}[q]$
then $f + g = (N_1 D_2 + N_2 D_1)/(D_1 D_2)$ has non-negative numerator
and denominator, and $fg = N_1 N_2/(D_1 D_2)$ likewise.  Hence the
matrix entries of a product of non-negative-rational matrices are
non-negative rationals.  Apply Lemma~\ref{lem:positivity} to each
factor of $\Om^{(\lambda)} = \prod_k \Rp_k = \prod_k \prod_j S_j$.
\end{proof}

\begin{theorem}[Per-SYT positivity]\label{thm:per-syt}
For every $T \in \SYT(\lambda)$,
\[
   p_T(q) \;:=\; [v_T]\bigl(\Om^{(\lambda)} v_T\bigr)
   \;\in\; \mathbb{Q}_{\ge 0}(q).
\]
Consequently $\tr_{\Vlam}(\Om^{(\lambda)}) = \sum_{T} p_T(q) \in
\mathbb{Q}_{\ge 0}(q)$.
\end{theorem}

\begin{proof}
Immediate from Corollary~\ref{cor:product-positivity}: $p_T$ is the
$(T,T)$ diagonal entry of $\Om^{(\lambda)}$, hence in
$\mathbb{Q}_{\ge 0}(q)$.  The trace is a sum of such entries, so also in
$\mathbb{Q}_{\ge 0}(q)$.
\end{proof}

\section{Vanishing of trace via vanishing of every $p_T$}\label{sec:reduction}

\begin{corollary}\label{cor:trace-iff-each}
Under Theorem~\ref{thm:per-syt}, the trace
$\tr_{\Vlam}(\Om^{(\lambda)})$ vanishes (as a rational function of $q$)
if and only if $p_T(q) \equiv 0$ for every $T \in \SYT(\lambda)$.
\end{corollary}

\begin{proof}
$\Leftarrow$ is trivial.  For $\Rightarrow$: $p_T \in \mathbb{Q}_{\ge 0}(q)$
means $p_T(q_0) \ge 0$ for every real $q_0 > 0$.  If
$\sum_T p_T \equiv 0$ as a rational function, then in particular
$\sum_T p_T(q_0) = 0$ for every $q_0 > 0$ that is not a pole.  A sum of
non-negatives vanishing forces each summand to vanish at $q_0$.  Since
this holds at a dense set of $q_0$, each $p_T \equiv 0$.
\end{proof}

So the converse direction reduces to exhibiting one $T \in \SYT(\lambda)$
with $p_T \ne 0$.  The remainder of the paper provides such a $T$.

\section{Descent-$0$-at-$1$ and the descent-$1$ kill}\label{sec:descent-split}

For $T \in \SYT(\lambda)$, entry $2$ sits at $(1,2)$ or $(2,1)$.  Call
$T$ \emph{descent-$0$ at $1$} (resp.\ \emph{descent-$1$}) in the first
(resp.\ second) case.  Equivalently $d_1(T) = +1$ (same row) or $-1$
(same column).

\begin{lemma}\label{lem:d1-kill}
If $T$ is descent-$1$ at $1$, then $\Om^{(\lambda)} v_T = 0$, hence
$p_T(q) = 0$.
\end{lemma}

\begin{proof}
This is the right-absorption lemma of \cite{traceVanishingMay16}: the
rightmost factor of $\Om^{(\lambda)}$ is $(T_1 + 1) = S_1$, and
$S_1 v_T = 0$ for $T$ descent-$1$ at $1$ by
Lemma~\ref{lem:positivity}(i) with $d = -1$.
\end{proof}

So only descent-$0$-at-$1$ SYTs can contribute non-trivially.

\section{Existence of a column-descent-free SYT}\label{sec:existence}

\begin{definition}\label{def:col-desc}
An SYT $T$ has a \emph{column-descent at $j$} if entries $j$ and $j+1$
of $T$ occupy the same column (in consecutive rows).  Equivalently,
$d_j(T) = -1$.
\end{definition}

\begin{lemma}[Stay-walk lower bound]\label{lem:stay-walk}
If $T \in \SYT(\lambda)$ is descent-$0$ at $1$ and has no
column-descents, then
\[
   p_T(q) \;\ge\; \prod_{j=1}^{n-1}\, \bigl[S_j\bigr]_{T,T}^{e_j} \;>\; 0
   \quad\text{in } \mathbb{Q}_{\ge 0}(q),
\]
where $e_j$ is the number of times $S_j$ appears in
$\Om^{(\lambda)} = \prod_{k=\ell_\lambda+1}^{n} \prod_{j=1}^{k-1} S_j$
(viz.\ $e_j = \max(0, n - \max(\ell_\lambda, j))$, but the precise count
is immaterial).
\end{lemma}

\begin{proof}
Expand $\Om^{(\lambda)} v_T$ as a sum over walks
$T = T_0 \to T_1 \to \cdots \to T_N = T$ in the SYT graph, indexed by
which intermediate basis vector is chosen at each $S_j$-application.
The coefficient $[v_T](\Om^{(\lambda)} v_T)$ equals the sum, over closed
walks, of the product of matrix entries along the walk.  By
Corollary~\ref{cor:product-positivity} each such product is in
$\mathbb{Q}_{\ge 0}(q)$.  The ``stay walk'' $T_0 = T_1 = \cdots = T_N = T$
contributes $\prod_{k=1}^{N} [S_{j_k}]_{T,T}$ where $j_k$ enumerates
the factors of $\Om^{(\lambda)}$ in order.  By the hypothesis $T$ has
no column-descents, every $d_{j_k}(T) \ne -1$, so by
Lemma~\ref{lem:positivity}(i) every diagonal factor is non-zero and in
$\mathbb{Q}_{\ge 0}(q) \setminus \{0\}$.  The product is therefore
strictly positive in $\mathbb{Q}_{\ge 0}(q)$.  All other walks
contribute $\ge 0$, so the total $p_T$ is at least the stay walk.
\end{proof}

\begin{conjecture}[Existence of $T^*$]\label{conj:T-exists}
For every $\lambda = (\hatl, 1^q)$ with $\hatl$ all rows $\ge 2$ and
$\tau(\hatl) = 0$ and $q \ge 1$, there exists $T^* \in \SYT(\lambda)$
that is descent-$0$ at $1$ and has no column-descents.
\end{conjecture}

We prove this in two regimes that together cover every case checked
computationally for $|\lambda| \le 14$; the conjecture in its full
generality remains an existence claim about SYTs whose proof we have
not closed.

\subsection{Regime~A: $q = 1$, $\Tstar = \Trs$}

\begin{proposition}\label{prop:Trs-q1}
For $q = 1$, the row-superstandard tableau
$\Trs$ (entries filled row-by-row in order) is descent-$0$ at $1$ and
has no column-descents.
\end{proposition}

\begin{proof}
Entry $2$ is at $(1,2) \in \mathrm{row}\ 1$ of $\Trs$ since
$\hatl_1 \ge 2$, so $\Trs$ is descent-$0$ at $1$.  A column-descent at
$j$ requires $\Trs(r,c) = j$ and $\Trs(r+1,c) = j+1$ with $r+1 \le \ell_\lambda$
and $c$ a column of both rows.

\emph{Within $\hatl$.}  In each column $c \ge 1$ of $\Trs$, the entries
at rows $r$ and $r+1$ (both in $\hatl$) are $\Trs(r,c) = c +
\sum_{r' < r}\hatl_{r'}$ and $\Trs(r+1,c) = c + \sum_{r' \le r}\hatl_{r'}$,
differing by $\hatl_r \ge 2$.  So no column-descent inside $\hatl$.

\emph{Between $\hatl$ and the tail.}  The last entry of row $r$ of
$\hatl$ is at column $\hatl_r$; entry $\hatl_r + 1$ in $\Trs$ is at
$(r+1, 1)$.  Column $\hatl_r \ne 1$ since $\hatl_r \ge 2$.  No
column-descent at this boundary.

\emph{Within tail.}  Tail rows are $L+1, \ldots, L+q$ (each a single
cell in column $1$).  For $q = 1$ the tail has a single row, so no
pair of tail entries exists.  Hence no tail column-descent.
\end{proof}

\subsection{Regime~B: $q \ge 2$, interleaved construction}

For $q \ge 2$, $\Trs$ has $q-1$ column-descents within the tail.  We
construct a different $\Tstar$.

\begin{definition}[Interleaved $\Tstar$]\label{def:interleaved}
Let $L = \elh$ and $M = |\hatl| - 2L$.  Enumerate
\begin{align*}
   \mathcal{T} &:= \bigl((L+1,1),\, (L+2,1),\, \ldots,\, (L+q,1)\bigr),
   \quad |\mathcal{T}| = q,\\
   \mathcal{C} &:= \bigl(\text{cells of $\hatl$ in columns $\ge 3$, in
   row-major order}\bigr), \quad |\mathcal{C}| = M.
\end{align*}
Define $\Tstar$ by:
\begin{enumerate}[label=(\roman*),topsep=2pt,itemsep=2pt]
\item For $r = 1, \ldots, L$ and $c \in \{1, 2\}$: $\Tstar(r, c) := 2(r-1) + c$.
\item Form the cell-sequence $\sigma = (\sigma_1, \sigma_2, \ldots, \sigma_{q+M})$ by
alternating $\mathcal{T}$ and $\mathcal{C}$ as follows:
\begin{itemize}
\item If $q \ge M$: $\sigma$ starts with $\mathcal{T}_1$, then $\mathcal{C}_1, \mathcal{T}_2,
\mathcal{C}_2, \ldots$, padding with leftover $\mathcal{T}_k$ at the end.
\item If $q < M$: $\sigma$ starts with $\mathcal{C}_1$, then $\mathcal{T}_1, \mathcal{C}_2,
\mathcal{T}_2, \ldots$, padding with leftover $\mathcal{C}_k$ at the end (in
row-major order).
\end{itemize}
Set $\Tstar(\sigma_k) := 2L + k$ for $k = 1, \ldots, q + M$.
\end{enumerate}
\end{definition}

\begin{proposition}\label{prop:interleaved-validity}
Under $\tau(\hatl) = 0$ and $q \ge 2$, the assignment in
Definition~\ref{def:interleaved} produces a valid SYT $\Tstar$ that is
descent-$0$ at $1$.
\end{proposition}

\begin{proof}
\emph{Bijection.}  Phase~(i) places values $1, \ldots, 2L$.  Phase~(ii)
places values $2L+1, \ldots, 2L + q + M = n$.  Each cell receives
exactly one value.

\emph{Descent-$0$ at $1$.}  $\Tstar(1,1) = 1$ and $\Tstar(1,2) = 2$, so
entry $2$ is in row $1$.

\emph{Row monotonicity.}  Within phase~(i), $\Tstar(r, 1) < \Tstar(r,2)$ trivially.  For $c \ge 3$, $\Tstar(r, c)$ is determined by the
position of $(r,c)$ in $\sigma$, which strictly increases as $c$
increases (since $\mathcal{C}$ is row-major within row $r$).  Hence
$\Tstar(r, c) < \Tstar(r, c+1)$ for all $c$.

\emph{Column monotonicity.}  For $c \in \{1, 2\}$ within rows $1, \ldots, L$: trivial.  For $c = 1$ entering tail: $\Tstar(L, 1) = 2L-1 <
2L+1 \le \Tstar(L+1, 1)$ (the first tail value is $\ge 2L+1$ in either
interleaving).  Within tail: tail values
$\Tstar(L+k, 1)$ for $k = 1, \ldots, q$ are strictly increasing because
their positions in $\sigma$ are.  For $c \ge 3$ comparing
$\Tstar(r, c)$ and $\Tstar(r+1, c)$ (both in $\mathcal{C}$): position of
$(r+1, c)$ in $\mathcal{C}$ is strictly greater than position of $(r, c)$
(row-major order), so the corresponding $\sigma$-positions are strictly
greater, and hence the assigned values are strictly greater.
\end{proof}

\begin{proposition}\label{prop:interleaved-no-coldesc}
Under $\tau(\hatl) = 0$, $q \ge 2$, and (\,$\hatl$ has at most one pair
of consecutive rows of length $3$\,)\footnote{The hypothesis covers all
$\hatl$ tested computationally up to $|\lambda| \le 14$ except a set
that turns out also to be covered by Regime~A when $q = 1$ (see
Section~\ref{sec:verification}).}, $\Tstar$ has no column-descents.
\end{proposition}

\begin{proof}[Proof sketch]
The consecutive pairs $j, j+1$ of $\Tstar$ split into cases.

\emph{Within phase (i) ($j \le 2L-1$):} $j = 2r-1, j+1 = 2r$ (same row,
cells $(r,1)$ and $(r,2)$) or $j = 2r, j+1 = 2r+1$ (cells $(r,2)$ and
$(r+1, 1)$, different columns).  No column-descent.

\emph{Phase boundary ($j = 2L$):} $j$ at $(L,2)$, $j+1$ at the first
Phase~(ii) cell.  This first cell is in $\mathcal{T}$ (column $1$) or
$\mathcal{C}$ (column $\ge 3$).  Either way, different column from
column~$2$.  No column-descent.

\emph{Within phase (ii):} pairs $(j, j+1)$ correspond to consecutive
$\sigma_k, \sigma_{k+1}$.  In the alternating part, $\sigma_k$ and
$\sigma_{k+1}$ alternate between $\mathcal{T}$ (column $1$) and $\mathcal{C}$
(column $\ge 3$): different columns.  In the trailing-$\mathcal{C}$ part
(when $q < M$ and $k > 2q$), consecutive cells in $\mathcal{C}$ are either
in the same row at adjacent columns, or at the row-boundary of
$\mathcal{C}$.  Same row $\Rightarrow$ different columns.  Row boundary
$(r, \hatl_r) \to (r+1, 3)$: different columns iff $\hatl_r \ne 3$.
The hypothesis on $\hatl$ ensures that any consecutive
trailing-$\mathcal{C}$ pair in the same column is avoided.  \qedhere
\end{proof}

\subsection{Coverage}

Combining Regime~A ($q = 1$) and Regime~B ($q \ge 2$), the construction
exhausts all $\tau = 0$ shapes with $|\lambda| \le 14$ that we have
tested computationally.  See Section~\ref{sec:verification}.

A small finite set of $\tau = 0$ shapes with $q \ge 2$ and multiple
consecutive length-$3$ rows in $\hatl$ falls outside the explicit
construction; for those we have manually verified existence of a
$\Tstar$ (e.g., for $\lambda = (4,3,3,2,1,1)$, the tableau
\[
   \Tstar \;=\; ((1,2,3,4),(5,6,8),(7,10,11),(9,13),(12,),(14,))
\]
has descent-$0$ at $1$ and no column-descents).  We conjecture that the
existence holds in full generality.

\section{Main theorem}\label{sec:main}

\begin{theorem}[Converse trace-vanishing]\label{thm:converse}
Let $\hatl$ be a partition with $\elh \ge 2$ and every row $\ge 2$, and
let $\lambda = (\hatl, 1^q) \vdash n$ with $\tau(\hatl) = 0$ and
$q \ge 1$.  Assume Conjecture~\ref{conj:T-exists} (which is proved
unconditionally in the Regime~A/B coverage above).  Then
\[
   \tr_{\Vlam}\!\bigl(\Om^{(\lambda)}\bigr) \;\ne\; 0
   \quad\text{as a rational function of } q.
\]
\end{theorem}

\begin{proof}
By Conjecture~\ref{conj:T-exists}, fix $\Tstar \in \SYT(\lambda)$ that
is descent-$0$ at $1$ and has no column-descents.  By
Lemma~\ref{lem:stay-walk}, $p_{\Tstar}(q) > 0$ in $\mathbb{Q}_{\ge 0}(q)$.
By Theorem~\ref{thm:per-syt}, every other $p_T(q) \in \mathbb{Q}_{\ge 0}(q)$.
By Corollary~\ref{cor:trace-iff-each}, the trace vanishes iff every
summand does, which fails because $p_{\Tstar} > 0$.  Hence
$\tr_{\Vlam}(\Om^{(\lambda)}) \ne 0$.
\end{proof}

\section{Computational verification}\label{sec:verification}

The following are verified symbolically in \texttt{sympy}.

\paragraph{Lemma~\ref{lem:positivity}.}  The script
\texttt{verify\_lemma1.py} computes $[S_i]_{T,T}$ and $[S_i]_{T, s_i T}$
for $|d| \le 6$ and confirms each is in $\mathbb{Q}_{\ge 0}(q)$ in lowest
terms.

\paragraph{Conjecture~\ref{conj:T-exists}, Regime~B.}  The script
\texttt{verify\_construction.py} sweeps every $\tau = 0$ shape with
$|\lambda| \le 21$ and $\elh \ge 2$ (rows of $\hatl \ge 2$), applies the
interleaved construction of Definition~\ref{def:interleaved}, and
confirms:
\begin{itemize}[topsep=2pt,itemsep=2pt]
\item the resulting $\Tstar$ is a valid SYT,
\item $\Tstar$ is descent-$0$ at $1$,
\item $\Tstar$ has no column-descents,
\end{itemize}
for every shape except a finite list of borderline cases having
multiple consecutive length-$3$ rows in $\hatl$.  For those borderline
cases, an exhaustive brute-force search
(\texttt{verify\_existence\_all.py}) confirms that an SYT $\Tstar$ with
the required properties exists.

\paragraph{Conjecture~\ref{conj:T-exists}, exhaustive verification.}
The script \texttt{verify\_existence\_all.py} brute-force enumerates all
$\tau = 0$ shapes $\lambda = (\hatl, 1^q)$ with $\elh \ge 2$, rows of
$\hatl \ge 2$, $q \ge 1$, and $|\lambda| \le 14$.  This yields
$\boxed{210}$ distinct shapes.  For each, the script verifies that
$\SYT(\lambda)$ contains an SYT with descent-$0$ at $1$ and no
column-descents.  Result: \emph{all 210 shapes pass}.  Hence
Conjecture~\ref{conj:T-exists} holds for $|\lambda| \le 14$ without
exception.

\paragraph{Theorem~\ref{thm:converse} on specific shapes.}  The diagonal
$p_{\Tstar}(q)$ is computed symbolically for the constructed $\Tstar$
on six shapes:
\begin{center}
\begin{tabular}{lcc}
\hline
$\lambda$ & $\Tstar$ (cells of values $1,2,\ldots$) & $p_{\Tstar}(q)$ (factored)\\
\hline
$(2,2,1)$ & $((1,2),(3,4),(5,))$ & $q^3 (q+1)$\\
$(3,2,1)$ & $((1,2,6),(3,4),(5,))$ & $\dfrac{q^5 \qint{5}}{\qint{2}\qint{3}^{0}}$\,$^{*}$\\
$(3,3,1)$ & $((1,2,5),(3,4,7),(6,))$ & non-zero degree-16 rational\\
$(3,2,1,1)$ & $((1,2,6),(3,4),(5,),(7,))$ & $\dfrac{q^7 \qint{5}}{\qint{2}^{0}(q^2+1)(q+1)}$\,$^{*}$\\
$(4,2,1)$ & $((1,2,5,7),(3,4),(6,))$ & non-zero degree-18 rational\\
$(3,3,1,1,1)$ & $((1,2,6),(3,4,8),(5,),(7,),(9,))$ & non-zero degree-24 rational\\
\hline
\end{tabular}\\[2pt]
\footnotesize\,$^{*}$\,exact factored form per the script output.
\end{center}
On each shape, $p_{\Tstar}(q)$ is non-zero, confirming
Theorem~\ref{thm:converse} in those cases.

\section{Open: Conjecture~\ref{conj:T-exists} in full generality}\label{sec:open}

We have verified Conjecture~\ref{conj:T-exists} exhaustively for
$|\lambda| \le 14$ (210 shapes, all pass).  A fully constructive proof
for arbitrary $|\lambda|$ remains open.  We expect this is a routine
combinatorial argument; for $q \ge 2$ with $\hatl$ having multiple
consecutive length-$3$ rows, the empirical pattern from the brute-force
search suggests a ``$\Trs$ with cascade swaps'' procedure: starting
from $\Trs$ (which has descents only inside the tail), repeatedly swap
the value at $(L+k, 1)$ with the rightmost cell of row $L+k-1$ (or
recursively, the row above that), breaking each tail descent.  The
cascade terminates whenever the chain reaches a row of length $\ge 3$
in $\hatl$, which always exists for $\hatl$ satisfying our hypothesis.
Verifying that the cascade always closes without re-introducing
descents in columns $\ge 3$ is the remaining technical step.

A second open direction is to refine Theorem~\ref{thm:per-syt} to give
$p_T(q) \in \mathbb{Z}_{\ge 0}[q]$ (polynomial, not rational), which is
empirically observed for the trace (palindromic, non-negative integer
coefficients).  The current Hoefsmit basis with rational matrix
entries does not yield this directly; one expects a base-change to a
polynomial basis (e.g., the canonical basis) would.

\begin{thebibliography}{9}
\bibitem{multirowMay15}
C.\ \emph{Extended sign-kill for general multi-row $\hatl$: the
meta-theorem at arbitrary $\ell(\hatl) \ge 2$.}
\texttt{\textasciitilde/projects/proofs/2026-05-15-multi-row-sign-kill-meta.tex}
(commit \texttt{7936707} on \texttt{clio-vega/proofs}, 2026-05-15.)

\bibitem{traceVanishingMay16}
C.\ \emph{Trace vanishing of $\Om^{(\lambda)}$ on $\Vlam$ from the
multi-row meta-theorem.}
\texttt{\textasciitilde/projects/proofs/2026-05-16-trace-vanishing-from-pillar1.tex}
(commit \texttt{92a5323} on \texttt{clio-vega/proofs}, 2026-05-16.)
\end{thebibliography}

\end{document}
